%======================================================================
\section{More about $\Gamma$}\label{sec:more}
%======================================================================

Theorem~A asserts that $\Gamma$ is finitely presented, and the introduction claims that $\Gamma$ embeds in Brin's group $3V$. This section
proves both, and then shows that the F\o lner function of $\Gamma$ is doubly exponential, so that the hash models of Section~\ref{sec:profile}
are exponentially smaller than any model built from a F\o lner set. Nothing here is used in the profile bounds of
Theorem~\ref{thm:main-group}, parts (a) to (c), or in any proof of Part~II, where the embedding in $3V$ serves only as a second
check of the Choi rank; part~(d) of that theorem, finite presentation, is
Proposition~\ref{prop:fp}. A reader heading for Theorem~B can skip this section and Section~\ref{sec:ladder}.

\subsection{A finite presentation}\label{sec:fp}

The proof that $\Gamma$ is finitely presented has three steps. A presentation lemma for groups of the form $F(\mathcal X)\rtimes H$ reduces
everything to finitely many $H$-orbits of relations. Steinberg's relations for the transvections, with Johnson's presentation of $H$ and the
relations of the translations, then present the affine group $A$, and omitting the translations presents $L$. Cornulier's criterion for
permutational wreath products presents $\Gamma$, because $A$ acts doubly transitively on the configurations with finitely generated
stabilizers.

The certificate of Section~\ref{sec:certificate} is a finite set of relations, but it does not present $\Gamma$. The group it presents has
abelianization $\Z^2\times\Z/2\times\Z/4$, in which $c$ has order four; this is a computation with exponent sums, for instance $a^2$, $b^2=(ac^4)^2$
and $(ab)^3$ give $2a=8c=12c=0$. In $\Gamma$, by contrast, $c=d\,ab$ is a product of commutators, since $d=[E_{2\leftarrow3},E_{3\leftarrow1}]$
and $ab$ is a three-cycle.

\begin{proposition}[Finite presentation]\label{prop:fp}
The groups $L$, $A$ and $\Gamma$ are finitely presented.
\end{proposition}

\begin{proof}
\emph{Stabilizers.} For a finite set $D\subset Y$ the pointwise stabilizer $H_{(D)}$ of $D$ in $H$ is isomorphic to $H$. Renumbering the points of
each ray outside $D$ in increasing order turns $Y\setminus D$ into three rays, and a permutation fixing $D$ is eventually a translation on the
old rays exactly when it is on the new ones. So $H_{(D)}$ is finitely generated. Since $H$ contains $\FSym(Y)$, it acts transitively on ordered
$k$-tuples of distinct points for every $k$.

\emph{A presentation lemma.} Let a finitely presented group $H$ act on a set $\mathcal X$ with finitely many orbits and finitely generated
stabilizers, and let $F(\mathcal X)$ be the free group on $\mathcal X$. Then $F(\mathcal X)\rtimes H$, with $H$ permuting the free generators, is
finitely presented. Take a presentation of $H$, adjoin one representative $x_i$ of each orbit, and add the relations $[x_i,k]=1$ for $k$ in a
finite generating set of the stabilizer of $x_i$. In the group so presented, $hx_ih^{-1}$ depends only on $h\cdot x_i$, and the group is
isomorphic to $F(\mathcal X)\rtimes H$: both are the iterated amalgam of $H$ with the groups $\mathrm{Stab}(x_i)\times\langle x_i\rangle$ over
the stabilizers $\mathrm{Stab}(x_i)$. Imposing in addition a set of relations that is a finite union of $H$-orbits keeps the group finitely
presented, since one representative of each orbit suffices.

\emph{A presentation of $A$.} For distinct $i,j\in Y$ let $x_{ij}$ be the transvection $e_j\mapsto e_j+e_i$, let $t_i=\tau_{e_i}$, and let
$\varsigma_{ij}\in\FSym(Y)\le H$ be the transposition of $i$ and $j$. Let $\widetilde A$ be the quotient of $F(\mathcal X)\rtimes H$, where
$\mathcal X=\{x_{ij}\}\sqcup\{t_i\}$ and $H$ relabels indices, by the relations
\begin{gather*}
x_{ij}^2=1,\qquad [x_{ij},x_{kl}]=1\ (i\ne l,\ j\ne k),\qquad [x_{ij},x_{jk}]=x_{ik}\ (i,j,k\text{ distinct}),\\
\varsigma_{ij}=x_{ij}x_{ji}x_{ij},\qquad t_i^2=1,\qquad [t_i,t_j]=1,\\
x_{ij}t_jx_{ij}^{-1}=t_it_j,\qquad x_{ij}t_kx_{ij}^{-1}=t_k\ (k\ne j).
\end{gather*}
They hold in $A$, so there is a surjection $\widetilde A\to A$. Each relation involves at most four points, so by transitivity the relations form
finitely many $H$-orbits; the stabilizers of $x_{ij}$ and $t_i$ are $H_{(\{i,j\})}$ and $H_{(\{i\})}$. Since $H$ is finitely
presented~\cite{Brown1987}, for instance by Johnson's presentation~\eqref{eq:johnson}~\cite{Johnson1999}, $\widetilde A$ is finitely presented, and
it remains to show that $\widetilde A\to A$ is injective. The linear relations are Steinberg's~\cite{Milnor1971}; we give the short direct argument.

Fix a finite set $D\subset Y$, number its points $1,\dots,m$ in an order unrelated to heights, and let $U_D\le\widetilde A$ be generated by the
$x_{ij}$ with $i<j$ in $D$. By the first three relations, $x_{1m},\dots,x_{m-1,m}$ are commuting involutions, and they generate a subgroup
normalized by the $x_{ij}$ with $i<j<m$. By induction $|U_D|\le2^{m-1}2^{(m-1)(m-2)/2}=2^{m(m-1)/2}$, the order of the upper unitriangular
group $\mathrm{UT}_m(\F_2)$, onto which $U_D$ maps. So $U_D\cong\mathrm{UT}_m(\F_2)$. Let $W_D=\Sym(D)\le H$ and $\varsigma_i=\varsigma_{i,i+1}$. By
the relation $\varsigma_{ij}=x_{ij}x_{ji}x_{ij}$ and relabeling, the subgroup $G_D$ generated by all $x_{ij}$ with $i,j\in D$ contains $W_D$ and is
generated by $U_D$ and the $\varsigma_i$. We claim $G_D=U_DW_DU_D$; since this set contains $1$, it suffices that it is closed under right
multiplication by $U_D$ and by each $\varsigma_i$. Every $b\in U_D$ is $x_{i,i+1}^\epsilon b_0$ with $b_0$ in the subgroup generated by the other
positive roots, the kernel of the entry $(i,i+1)$, which $\varsigma_i$ normalizes. So it suffices to treat $wx_{i,i+1}\varsigma_i$ with $w\in W_D$;
put $k=w(i)$ and $l=w(i+1)$. If $k<l$, then $wx_{i,i+1}\varsigma_i=x_{kl}w\varsigma_i\in U_DW_D$. If $k>l$, then $w\varsigma_i=\varsigma_{kl}w$ and
$x_{kl}\varsigma_{kl}=x_{lk}x_{kl}$, so $wx_{i,i+1}\varsigma_i=x_{kl}\varsigma_{kl}w=x_{lk}wx_{i,i+1}\in U_DW_DU_D$. The map $G_D\to\GL(\F_2^D)$ is
onto, since transvections generate $\GL(\F_2^D)$. If $b_1wb_2$ maps to the identity of $\GL(\F_2^D)$, the permutation matrix of $w$ is upper
unitriangular, so $w=1$ in $H$, and then $b_1b_2=1$ in $U_D$. Hence $G_D\cong\GL(\F_2^D)$. Every element of the subgroup $\widetilde J$ generated
by all $x_{ij}$ lies in some $G_D$, so $\widetilde J\cong\GL_{\rm fin}(Y)$. By relabeling, every element of the subgroup generated by $\widetilde J$
and $H$ has the form $jh$. If it maps to $1$, then $h\in\GL_{\rm fin}(Y)\cap H=\FSym(Y)$, so $h$ is a product of transpositions, each equal in
$\widetilde A$ to a word in the $x_{ij}$, and $jh\in\widetilde J$ maps to $1$, hence is $1$. Finally the $t_i$ generate a normal elementary abelian
subgroup. It is a quotient of $\F_2^{(Y)}$ and maps onto $V$ by $t_i\mapsto e_i$, hence isomorphically, and the linear part acts on it as on $V$.
Hence every element of $\widetilde A$ is $t\ell$ with $t$ in that subgroup. If its image is the identity affine map, evaluating at $0$ gives
$t=1$, and then $\ell=1$. So $\widetilde A\cong A$. Omitting the $t_i$ gives $L$.

\emph{The lamps.} The group $A$ acts on $V$ doubly transitively: a translation carries $v$ to $0$, and a finitary linear map carries any nonzero
vector to $e_1$. The stabilizer of $0$ is $L$, which is generated by $p$, $q$ and $d$, because the $H$-conjugates of $d$ are all the transvections
$E_{y\leftarrow y'}$. So $A$ is finitely presented, its point stabilizers are finitely generated, and it has two orbits on $V\times V$. By
Cornulier's criterion for permutational wreath products~\cite[Theorem~1.1]{Cornulier2006}, $\Gamma=S_3\wr_VA$ is finitely presented. Explicitly,
let $\mathcal R_A$ present $A$ on generators that include $p,q,d,t$. Add the generators $a,b$ and the relations $a^2$, $b^2$, $(ab)^3$, $[a,\ell]$
and $[b,\ell]$ for $\ell\in\{p,q,d\}$, and $[u,tvt^{-1}]$ for $u,v\in\{a,b\}$. The result presents $\Gamma$. The lamp at $v\in V$ is well defined
as a conjugate of $\langle a,b\rangle$, because $L$ centralizes $\langle a,b\rangle$; and double transitivity carries the commutation of the lamps
at $0$ and at $e_1$ to every pair of lamps.
\end{proof}

\subsection{An embedding in a Brin--Thompson group}\label{sec:3V}

This subsection places $\Gamma$ inside a finitely presented simple group whose soficity is open, so that this group has a chunk with finite,
nearly exponential profile. Let $\mathfrak C=\{0,1\}^{\N}$ be Cantor space. For $n\ge1$, Brin's group $nV$~\cite{Brin2004} consists of the
homeomorphisms $h$ of $\mathfrak C^n$ given by a finite table of prefix replacements. Such a table is a pair of partitions of $\mathfrak C^n$ into
the same number of rectangles $u_1\mathfrak C\times\dots\times u_n\mathfrak C$, with $u_1,\dots,u_n$ finite binary words, matched in pairs. The
map $h$ sends each rectangle of the first partition onto its partner by replacing the $n$ prefixes and keeping the $n$ suffixes. Thus $1V$ is Thompson's group
$V$, and $nV$ embeds in $n'V$ for $n\le n'$ by acting trivially on the additional coordinates. The groups $nV$ are finitely presented, by Brin for
$2V$~\cite{Brin2005} and by Hennig and Matucci for every $n$~\cite{HennigMatucci2012}, and simple~\cite{Brin2010}. As for $V$, it is not known
whether they are sofic. The letter $V$ in these names has nothing to do with our space of configurations, which is also called $V$.

\begin{proposition}[Embedding in $3V$]\label{prop:3V}
$\Gamma$ embeds in $3V$, and hence in $nV$ for every $n\ge3$. Consequently $3V$ has a chunk $E_{3V}$ with
\[
\exp\Bigl(\frac{c\,r}{(\log r)^{12.75}}\Bigr)\ \le\ \sigma_{E_{3V}}(r)\ \le\ \exp(C\,r\log r)\qquad(r\ge2).
\]
\end{proposition}

\begin{proof}
\emph{Coding configurations.} For a finite set $S=\{j_1<\dots<j_m\}$ of positive integers, with $j_0=0$, let
\[
e(S)=\Bigl(\prod_{k=1}^m1\,0^{\,j_k-j_{k-1}-1}\,1\Bigr)\,0,
\]
so that $e(\emptyset)=0$, $e(\{1\})=110$ and $e(\{2\})=1010$. Read a word from the left: at a record boundary the letter $0$ ends the word and the
letter $1$ opens a record $10^k1$, which ends at its next letter $1$. So the words $e(S)$ form a prefix code. A finite word whose reading has not
ended can be extended to one that begins with some $e(S)$, by appending $0$ at a record boundary or $10$ inside a record. Hence the cylinders
$e(S)\mathfrak C$ are disjoint and their union is dense in $\mathfrak C$. For $v\in V$ and $i=1,2,3$ let $S_i(v)=\{j:v(i,j)=1\}$, let
\[
Z_v=e(S_1(v))\mathfrak C\times e(S_2(v))\mathfrak C\times e(S_3(v))\mathfrak C ,
\]
and let $\iota_v:\mathfrak C^3\to Z_v$ prefix the three codewords. The rectangles $Z_v$ are disjoint, and their union is dense in $\mathfrak C^3$.

\emph{Prepending a site.} Define $s_0,s_1:\mathfrak C\to\mathfrak C$ by $s_0(0\xi)=0\xi$, $s_0(1\xi)=10\xi$ and $s_1(\xi)=11\xi$. Their images
$0\mathfrak C\cup10\mathfrak C$ and $11\mathfrak C$ partition $\mathfrak C$, and, with $S+1=\{j+1:j\in S\}$,
\begin{equation}\label{eq:prepend}
s_0\bigl(e(S)\xi\bigr)=e(S+1)\,\xi,\qquad s_1\bigl(e(S)\xi\bigr)=e\bigl(\{1\}\cup(S+1)\bigr)\,\xi :
\end{equation}
an empty site in front lengthens the first gap by one, and an occupied site in front adds the record $11$.

\emph{The affine group.} For $b\in\{0,1\}$ put
\begin{gather*}
\hat P(x,s_b(y),z)=(s_b(x),y,z),\qquad \hat Q(x,y,s_b(z))=(s_b(x),y,z),\\
\hat T(s_b(x),y,z)=(s_{1-b}(x),y,z),
\end{gather*}
and let $\hat D$ fix $s_0(\mathfrak C)\times\mathfrak C^2$ pointwise and send $(11x,y,z)$ to $(11x',y,z)$, where $\hat T(x,y,z)=(x',y,z)$. These maps lie
in $3V$. In prefix tables, with $\epsilon$ the empty word, $\hat P$ replaces the prefixes $(\epsilon,11,\epsilon)$, $(0,0,\epsilon)$, $(1,0,\epsilon)$,
$(0,10,\epsilon)$ and $(1,10,\epsilon)$ by $(11,\epsilon,\epsilon)$, $(0,0,\epsilon)$, $(10,0,\epsilon)$, $(0,1,\epsilon)$ and $(10,1,\epsilon)$. $\hat Q$ does
the same with the second and third coordinates exchanged. $\hat T$ exchanges the first-coordinate prefixes $0\leftrightarrow110$ and
$10\leftrightarrow111$. $\hat D$ exchanges $110\leftrightarrow11110$ and $1110\leftrightarrow11111$ and keeps $0$ and $10$. The permutation $p$ moves
the first point of ray~$2$ to the first point of ray~$1$ and shifts the other points of these rays, and $q$ does the same with ray~$3$. The element
$t$ changes the value of a configuration at the point $1$, and $d$ adds the value at $1$ to the value at $2$. So~\eqref{eq:prepend} gives, for
every $v\in V$,
\[
\hat P\iota_v=\iota_{p(v)},\qquad \hat Q\iota_v=\iota_{q(v)},\qquad \hat T\iota_v=\iota_{t(v)},\qquad \hat D\iota_v=\iota_{d(v)} .
\]
By the proof of Lemma~\ref{lem:structure}(b), $p,q,t,d$ generate $A$. If a word in them represents $g\in A$, the same word in
$\hat P,\hat Q,\hat T,\hat D$ maps $\iota_v(\xi)$ to $\iota_{g(v)}(\xi)$ for all $v$ and $\xi$. If $g=1$, it is the identity on the dense set
$\bigcup_vZ_v$, hence the identity. So $p,q,t,d\mapsto\hat P,\hat Q,\hat T,\hat D$ defines a homomorphism $\rho:A\to3V$ with
\begin{equation}\label{eq:3V-equivariant}
\rho(g)\,\iota_v=\iota_{g(v)}\qquad(g\in A,\ v\in V),
\end{equation}
and $\rho$ is injective: if $g\ne1$, then $g(v)\ne v$ for some $v$, and $\rho(g)$ maps $Z_v$ onto the disjoint rectangle $Z_{g(v)}$.

\emph{The lamps.} Let $\beta_1=0$, $\beta_2=10$ and $\beta_3=11$, and let $S_3$ act faithfully on $\mathfrak C^3$ by the prefix replacements
$\theta(\lambda)(\beta_i\xi,y,z)=(\beta_{\lambda(i)}\xi,y,z)$. For $\lambda\in S_3$ and $v\in V$ let $\rho(\lambda_v)$ be $\iota_v\theta(\lambda)\iota_v^{-1}$ on
$Z_v$ and the identity off $Z_v$; it lies in $3V$, since the complement of a rectangle is a finite union of rectangles. Copies at different
configurations have disjoint supports and commute, each copy is faithful, and by~\eqref{eq:3V-equivariant}
$\rho(g)\rho(\lambda_v)\rho(g)^{-1}=\rho(\lambda_{g(v)})$ for $g\in A$. These are the defining relations of $\Gamma$ as a semidirect product, so $\rho$
extends to a homomorphism $\rho:\Gamma\to3V$. Suppose that $\rho(\gamma)=1$, where $\gamma$ has lamp part $\lambda:V\to S_3$ and affine part $g$. The
lamps preserve every $Z_v$, so $\rho(g)$ does too, and $g(v)=v$ for all $v$; so $g=1$. Then $\rho(\gamma)$ restricts on $Z_v$ to
$\iota_v\theta(\lambda(v))\iota_v^{-1}$, so $\lambda(v)=1$ for all $v$. Hence $\rho$ is injective. Explicitly, $\rho(a)$ exchanges the rectangles
$00\mathfrak C\times0\mathfrak C\times0\mathfrak C$ and $010\mathfrak C\times0\mathfrak C\times0\mathfrak C$. Also $\rho(c)=\hat D\rho(ab)$, where
$\rho(ab)$ permutes the rectangles with prefixes $(00,0,0)$, $(010,0,0)$ and $(011,0,0)$ cyclically in this order and $\hat D$ is the identity on
$Z_0$.

\emph{The chunk.} The sofic profile of a chunk depends only on its identity element and its partial multiplication
(Section~\ref{sec:criterion}), and an injective homomorphism preserves both. So the chunk $E_{3V}=\rho(E)$, for the chunk $E$ of
Theorem~\ref{thm:main-group}, has $\sigma_{E_{3V}}=\sigma_E$, and parts (a) and (b) of that theorem give the bounds.
\end{proof}

By Section~\ref{sec:thompson}, Thompson's group $V$ has chunks with finite superpolynomial profile, from its Houghton subgroups, and a chunk with
profile at least $\exp(c\,r/\log^2r)$, which is finite if $F$ is amenable. Proposition~\ref{prop:3V} gives finite, nearly exponential profile in
$3V$ with no assumption. Whether $\Gamma$ embeds in $V$ or in $2V$ is open (Section~\ref{sec:open}).

\subsection{The F\o lner function}\label{sec:folner}

The hash models of Proposition~\ref{prop:upper} are far smaller than any F\o lner set of $\Gamma$, and this subsection measures the gap. For
$r\ge1$ let $\textup{Føl}(r)$ be the least size of a finite nonempty set $F\subset\Gamma$ with $|Fs\setminus F|\le|F|/r$ for
$s\in\{p,q,t,c,a\}$.

\begin{proposition}[The F\o lner function]\label{prop:folner-lower}
For $r\ge12$,
\[
2^{\,2^{\lfloor r/12\rfloor-1}}\ \le\ \textup{Føl}(r)\ \le\ 2^{\,2^{3r+14}} .
\]
Thus the Følner function of $\Gamma$ is doubly exponential, while by Theorem~\ref{thm:main-group}(b) the sofic profile of every chunk is
$\exp(O(r\log r))$.
\end{proposition}

\begin{proof}
\emph{Lower bound.} Let $F$ be such a set and $l=\lfloor r/12\rfloor\ge1$. For $v\in\{0,1\}^l$ the word
$T_v=t^{v_1}p\,t^{v_2}p\cdots p\,t^{v_l}p^{1-l}$ is freely equal to $\prod_{j=1}^lp^{j-1}t^{v_j}p^{1-j}$, so it represents the translation by
$\sum_jv_je_j$, and it has length less than $3l$. The elements $a_v=T_vaT_v^{-1}$ are the copies of $a$ at $2^l$ distinct configurations, so they
are commuting involutions and form a basis of an elementary abelian group $Q$ of order $2^{M}$, $M=2^l$. Each $a_v$ is a product of fewer than
$6l$ generators and inverses. Since $|Fxy\setminus F|\le|Fx\setminus F|+|Fy\setminus F|$ and $|Fx^{-1}\setminus F|=|Fx\setminus F|$, we get
$|Fa_v\setminus F|\le6l|F|/r\le|F|/2$, and so $\sum_v|Fa_v\setminus F|\le M|F|/2$.

On the other hand, for every finite $S\subset Q$,
\begin{equation}\label{eq:cube}
\sum_v|Sa_v\setminus S|\ \ge\ |S|\,(M-\log|S|).
\end{equation}
Indeed, the left side is $M|S|-2e(S)$, where $e(S)$ is the number of pairs $\{x,xa_v\}\subset S$, and $e(S)\le\frac12|S|\log|S|$ by the
edge-isoperimetric inequality of the cube~\cite{Harper1964}. We recall the proof, by induction on $M$. Split $S$ by the last coordinate into
$S_0,S_1$ of sizes $s_0\le s_1$. Then $e(S)\le e(S_0)+e(S_1)+s_0\le\frac12(s_0\log s_0+s_1\log s_1)+s_0$, and this is at most $\frac12s\log s$ with
$s=s_0+s_1$, because the binary entropy function $h$ satisfies $h(x)\ge2x$ for $0\le x\le\frac12$. The right multiplications by the $a_v$ preserve
each coset $xQ$; applying~\eqref{eq:cube} to the sets $x^{-1}(F\cap xQ)$ and using $|F\cap xQ|\le|F|$ gives
$\sum_v|Fa_v\setminus F|\ge|F|(M-\log|F|)$. Hence $\log|F|\ge M/2=2^{l-1}$.

\emph{Upper bound.} Let $m=\lceil(r-1)/2\rceil$, $R=2m+3$ and $F=\bigcup_{z\in B_m}K_Rg_z$, a disjoint union because $\pi(g_z)=z$. By
Lemma~\ref{lem:clock}, $g_zs\in K_Rg_{z+\pi(s)}$ for $z\in B_m$ and each generator $s$, so $|Fs\setminus F|\le|K_R|(2m+1)=|F|/(2m+1)\le|F|/r$. With
$D=|Y_R|=6m+9$ we have $|K_R|=6^{2^D}\cdot2^D\cdot|\GL_D(\F_2)|$, so $\log|F|\le2^D\log6+D+D^2+2\log(2m+1)\le2^{D+2}\le2^{3r+14}$.
\end{proof}

The group $\Gamma$ is not LEF. This follows from $H_3\le\Gamma$ in the standard way. $H_3$ is finitely presented and not residually finite,
because the finitary alternating group of an infinite set is simple. A finitely presented LEF group is residually
finite~\cite{VershikGordon1997}, and subgroups of LEF groups are LEF. The chunk $E$ shows it too. Its profile is unbounded by
Theorem~\ref{thm:main-group}(a), and a chunk with unbounded profile embeds in no finite group (Section~\ref{sec:criterion}).
